\chapter{DF-069 - Block-Authored DOGFOOD}\label{record-109}

\noindent{\footnotesize\color{BookMuted}\path{docs/design/DF-069-block-authored-dogfood.md}\par}\smallskip

Linked legend: \href{https://github.com/flyingrobots/bijou/blob/1853ad902f6bf219326409e7f743a5a1d6517dbc/docs/legends/DF-dogfood-field-guide.md}{DF - DOGFOOD Field Guide}

\section{Sponsor human}\label{record-109-sponsor-human}

DOGFOOD should prove that Bijou applications can be authored from Blocks, not only that Blocks can be previewed inside DOGFOOD.

\section{Sponsor agent}\label{record-109-sponsor-agent}

An agent should be able to inspect DOGFOOD and identify the semantic blocks that make up the product surface without reverse-engineering bespoke render helpers.

\section{Hill}\label{record-109-hill}

DOGFOOD becomes a block-authored application at the semantic surface boundary: title screen, shell regions, navigation, articles, settings, story previews, block previews, and inspector panes are declared as Blocks. Components remain the lower-level rendering vocabulary inside those Blocks.

\section{Core Rule}\label{record-109-core-rule}

Blocks own product semantics. Components own leaf rendering.

That means this work should create and consume Blocks for surfaces such as:

\begin{itemize}
\tightlist
\item
  \texttt{T\allowbreak{}i\allowbreak{}t\allowbreak{}l\allowbreak{}e\allowbreak{}S\allowbreak{}c\allowbreak{}r\allowbreak{}e\allowbreak{}e\allowbreak{}n\allowbreak{}B\allowbreak{}l\allowbreak{}o\allowbreak{}c\allowbreak{}k}
\item
  \texttt{N\allowbreak{}a\allowbreak{}v\allowbreak{}i\allowbreak{}g\allowbreak{}a\allowbreak{}t\allowbreak{}i\allowbreak{}o\allowbreak{}n\allowbreak{}L\allowbreak{}i\allowbreak{}s\allowbreak{}t\allowbreak{}B\allowbreak{}l\allowbreak{}o\allowbreak{}c\allowbreak{}k}
\item
  \texttt{D\allowbreak{}o\allowbreak{}c\allowbreak{}u\allowbreak{}m\allowbreak{}e\allowbreak{}n\allowbreak{}t\allowbreak{}a\allowbreak{}t\allowbreak{}i\allowbreak{}o\allowbreak{}n\allowbreak{}A\allowbreak{}r\allowbreak{}t\allowbreak{}i\allowbreak{}c\allowbreak{}l\allowbreak{}e\allowbreak{}B\allowbreak{}l\allowbreak{}o\allowbreak{}c\allowbreak{}k}
\item
  \texttt{S\allowbreak{}e\allowbreak{}t\allowbreak{}t\allowbreak{}i\allowbreak{}n\allowbreak{}g\allowbreak{}s\allowbreak{}M\allowbreak{}e\allowbreak{}n\allowbreak{}u\allowbreak{}B\allowbreak{}l\allowbreak{}o\allowbreak{}c\allowbreak{}k}
\item
  \texttt{S\allowbreak{}t\allowbreak{}o\allowbreak{}r\allowbreak{}y\allowbreak{}b\allowbreak{}o\allowbreak{}o\allowbreak{}k\allowbreak{}W\allowbreak{}o\allowbreak{}r\allowbreak{}k\allowbreak{}b\allowbreak{}e\allowbreak{}n\allowbreak{}c\allowbreak{}h\allowbreak{}B\allowbreak{}l\allowbreak{}o\allowbreak{}c\allowbreak{}k}
\item
  \texttt{B\allowbreak{}l\allowbreak{}o\allowbreak{}c\allowbreak{}k\allowbreak{}P\allowbreak{}r\allowbreak{}e\allowbreak{}v\allowbreak{}i\allowbreak{}e\allowbreak{}w\allowbreak{}B\allowbreak{}l\allowbreak{}o\allowbreak{}c\allowbreak{}k}
\item
  \texttt{G\allowbreak{}u\allowbreak{}i\allowbreak{}d\allowbreak{}e\allowbreak{}I\allowbreak{}n\allowbreak{}s\allowbreak{}p\allowbreak{}e\allowbreak{}c\allowbreak{}t\allowbreak{}o\allowbreak{}r\allowbreak{}B\allowbreak{}l\allowbreak{}o\allowbreak{}c\allowbreak{}k}
\end{itemize}

It should not wrap every \texttt{b\allowbreak{}o\allowbreak{}x\allowbreak{}S\allowbreak{}u\allowbreak{}r\allowbreak{}f\allowbreak{}a\allowbreak{}c\allowbreak{}e\allowbreak{}(\allowbreak{})}, \texttt{m\allowbreak{}a\allowbreak{}r\allowbreak{}k\allowbreak{}d\allowbreak{}o\allowbreak{}w\allowbreak{}n\allowbreak{}(\allowbreak{})}, or list row in a Block.

\section{Architecture}\label{record-109-architecture}

The target composition is:

% Presentation of original arrow notation. See flow-manifest.json.
\begin{center}
\begin{tikzpicture}[book flow node/.style={draw=BookHeading,line width=.65pt,fill=BookPanel,rounded corners=1.5pt,align=center,inner sep=5pt,minimum height=10mm,font=\fontsize{9}{11}\selectfont,execute at begin node={\hyphenpenalty=10000\exhyphenpenalty=10000}},book flow edge/.style={-{Latex[length=2mm]},draw=BookHeading,line width=.7pt}]
\node[book flow node,text width=96mm] (n0) at (0,0.0) {DOGFOOD app};
\node[book flow node,text width=96mm] (n1) at (0,-1.55) {App\allowbreak{}Frame /\allowbreak{} App\allowbreak{}Shell composition};
\node[book flow node,text width=96mm] (n2) at (0,-3.1) {page and region Blocks};
\node[book flow node,text width=96mm] (n3) at (0,-4.65) {components};
\node[book flow node,text width=96mm] (n4) at (0,-6.2) {surfaces};
\draw[book flow edge] (n0) -- (n1);
\draw[book flow edge] (n1) -- (n2);
\draw[book flow edge] (n2) -- (n3);
\draw[book flow edge] (n3) -- (n4);
\end{tikzpicture}
\end{center}


The existing unidirectional data-binding rule still applies:

% Presentation of original arrow notation. See flow-manifest.json.
\begin{center}
\begin{tikzpicture}[book flow node/.style={draw=BookHeading,line width=.65pt,fill=BookPanel,rounded corners=1.5pt,align=center,inner sep=5pt,minimum height=10mm,font=\fontsize{9}{11}\selectfont,execute at begin node={\hyphenpenalty=10000\exhyphenpenalty=10000}},book flow edge/.style={-{Latex[length=2mm]},draw=BookHeading,line width=.7pt}]
\node[book flow node,text width=96mm] (n0) at (0,0.0) {business logic /\allowbreak{} providers};
\node[book flow node,text width=96mm] (n1) at (0,-1.55) {immutable snapshots};
\node[book flow node,text width=96mm] (n2) at (0,-3.1) {binding frames};
\node[book flow node,text width=96mm] (n3) at (0,-4.65) {blocks and views render};
\node[book flow node,text width=96mm] (n4) at (0,-6.2) {command intents};
\node[book flow node,text width=96mm] (n5) at (0,-7.75) {business logic owns the next state};
\draw[book flow edge] (n0) -- (n1);
\draw[book flow edge] (n1) -- (n2);
\draw[book flow edge] (n2) -- (n3);
\draw[book flow edge] (n3) -- (n4);
\draw[book flow edge] (n4) -- (n5);
\end{tikzpicture}
\end{center}


Block-authored DOGFOOD must not introduce hidden provider registries, callback backchannels, render-time refresh hooks, mutable view stores, or direct provider reads from rendering code.

\section{Storybook Implication}\label{record-109-storybook-implication}

The standalone Storybook app should stop being a parallel bespoke TUI. It should use the same framed shell path as DOGFOOD and expose a \texttt{S\allowbreak{}t\allowbreak{}o\allowbreak{}r\allowbreak{}y\allowbreak{}b\allowbreak{}o\allowbreak{}o\allowbreak{}k\allowbreak{}W\allowbreak{}o\allowbreak{}r\allowbreak{}k\allowbreak{}b\allowbreak{}e\allowbreak{}n\allowbreak{}c\allowbreak{}h\allowbreak{}B\allowbreak{}l\allowbreak{}o\allowbreak{}c\allowbreak{}k} contract so story preview work exercises the product composition model.

\section{Non-goals}\label{record-109-non-goals}

\begin{itemize}
\tightlist
\item
  no full visual redesign of DOGFOOD
\item
  no attempt to convert every low-level component into a Block
\item
  no new provider lifecycle system
\item
  no new localization runtime
\item
  no hidden global block registry
\item
  no rendered AppShell policy invented outside existing contracts
\end{itemize}

\section{Accessibility / Assistive Posture}\label{record-109-accessibility-assistive-posture}

Block-authored DOGFOOD should make semantic regions easier to lower across interactive, static, pipe, and accessible modes. Blocks should preserve facts about navigation, content, settings, previews, and inspector state without requiring agents or assistive modes to parse terminal art.

\section{Localization / Directionality Posture}\label{record-109-localization-directionality-posture}

Blocks should receive localized labels and copy through DOGFOOD's localization port, not through direct catalog access. Block metadata may describe the surface, but runtime-visible labels remain localized app data.

\section{Agent Inspectability / Explainability Posture}\label{record-109-agent-inspectability-explainability-posture}

Agents should be able to ask which DOGFOOD Blocks are active, what data requirements they declare, which command intents they expose, and which surface they render into.

\section{Linked Invariants}\label{record-109-linked-invariants}

\begin{itemize}
\tightlist
\item
  \href{https://github.com/flyingrobots/bijou/blob/cbd70612140e7a450d764b3c7be1c6af878b99bc/docs/invariants/docs-are-the-demo.md}{Docs Are the Demo}
\item
  \href{https://github.com/flyingrobots/bijou/blob/cbd70612140e7a450d764b3c7be1c6af878b99bc/docs/invariants/runtime-truth-wins.md}{Runtime Truth Wins}
\item
  \href{https://github.com/flyingrobots/bijou/blob/cbd70612140e7a450d764b3c7be1c6af878b99bc/docs/invariants/commands-change-state-effects-do-not.md}{Commands Change State, Effects Do Not}
\item
  \href{https://github.com/flyingrobots/bijou/blob/cbd70612140e7a450d764b3c7be1c6af878b99bc/docs/invariants/tests-are-the-spec.md}{Tests Are the Spec}
\end{itemize}

\section{Implementation Outline}\label{record-109-implementation-outline}

\begin{enumerate}
\def\labelenumi{\arabic{enumi}.}
\tightlist
\item
  Update BEARING so the next product gravity is block-authored DOGFOOD.
\item
  Add DOGFOOD-local block registry helpers for semantic app surfaces.
\item
  Add \texttt{S\allowbreak{}t\allowbreak{}o\allowbreak{}r\allowbreak{}y\allowbreak{}b\allowbreak{}o\allowbreak{}o\allowbreak{}k\allowbreak{}W\allowbreak{}o\allowbreak{}r\allowbreak{}k\allowbreak{}b\allowbreak{}e\allowbreak{}n\allowbreak{}c\allowbreak{}h\allowbreak{}B\allowbreak{}l\allowbreak{}o\allowbreak{}c\allowbreak{}k} and move Storybook toward the framed shell.
\item
  Add DOGFOOD Blocks for title, navigation, documentation article, settings, block preview, and guide inspector surfaces.
\item
  Prove the block-authored registry without rendering every block during discovery.
\item
  Keep visual changes incremental and scoped.
\end{enumerate}

\section{Tests To Write First}\label{record-109-tests-to-write-first}

\begin{itemize}
\tightlist
\item
  behavior tests proving the BlockLab app is framed rather than bespoke
\item
  behavior tests proving DOGFOOD block registry discovery does not call render
\item
  behavior tests proving core DOGFOOD surfaces have Block definitions
\item
  behavior tests proving localization and command intent boundaries stay declarative
\end{itemize}

\section{Playback}\label{record-109-playback}

DF-069A lands the block-authored DOGFOOD contract layer:

\begin{itemize}
\tightlist
\item
  DOGFOOD has a local, branded \texttt{D\allowbreak{}o\allowbreak{}g\allowbreak{}f\allowbreak{}o\allowbreak{}o\allowbreak{}d\allowbreak{}B\allowbreak{}l\allowbreak{}o\allowbreak{}c\allowbreak{}k\allowbreak{}R\allowbreak{}e\allowbreak{}g\allowbreak{}i\allowbreak{}s\allowbreak{}t\allowbreak{}r\allowbreak{}y}.
\item
  Registry discovery rejects loose block-shaped objects.
\item
  Registry discovery records surface ownership without calling \texttt{r\allowbreak{}e\allowbreak{}n\allowbreak{}d\allowbreak{}e\allowbreak{}r\allowbreak{}(\allowbreak{})}.
\item
  BlockLab now has a \texttt{B\allowbreak{}l\allowbreak{}o\allowbreak{}c\allowbreak{}k\allowbreak{}L\allowbreak{}a\allowbreak{}b\allowbreak{}W\allowbreak{}o\allowbreak{}r\allowbreak{}k\allowbreak{}b\allowbreak{}e\allowbreak{}n\allowbreak{}c\allowbreak{}h\allowbreak{}B\allowbreak{}l\allowbreak{}o\allowbreak{}c\allowbreak{}k} contract and the interactive BlockLab entrypoint uses the AppFrame shell.
\item
  DOGFOOD surface blocks now cover title, navigation, documentation article, block preview, guide inspector, settings, and BlockLab workbench surfaces.
\end{itemize}

\section{Retrospective}\label{record-109-retrospective}

This slice deliberately stops at semantic surface contracts and the BlockLab frame migration. It does not rewrite every DOGFOOD renderer yet, does not move provider lifecycle policy, and does not convert leaf components into Blocks.

The final review gate also tightened the DOGFOOD i18n debt scanner so framed BlockLab machine tokens such as key names, event discriminants, mode literals, pane overflow policy, tones, and internal thrown errors do not count as visible English copy. BlockLab labels and rendered text still count as localization debt until they move behind a localization port.

